\section{Introduction}

\IEEEPARstart{C}{Computational} Flow Dynamics (CFD) simulation is becoming extremely important to speed up prototyping and validating new designs for large-scale models, which are difficult to develop by a purely analytical method. For instance, reservoir simulators are used extensively by reservoir engineers nowadays to analyze the flow of fluids and predict oil and gas production. This applies both to managing current and developing new fields. The core of such a CFD simulation is finding the solution to a set of partial differential equations (PDEs) constrained on some domain of interest. This is a two-step process, where the PDEs are first discretized over the defined domain using numerical techniques such as the finite-difference, -element, or -volume method, and second, solving the obtained linear and nonlinear systems using iterative linear solvers such as Krylov subspace solvers. However, solving the linear and nonlinear systems produced by these finite methods is usually time-consuming and complex because the linear systems are sparse and ill-conditioned. Furthermore, to improve the realism of the simulation, the number of cells required to be modeled is increasing at a steady pace. This leads to bigger systems, which in turn leads to higher run times. To accommodate these bigger systems and improve the scalability of the simulation, parallelization can be performed in two major directions: multi-core CPUs and many-core GPUs. In this work, we focus on the latter using the OPM project.

The Open Porous Media (OPM) \cite{OPM} project encourages open and reproducible research for modeling and simulation of porous media. OPM supports several types of reservoir simulation, including black oil. The project is split into several modules, e.g., opm-simulators, and is built on top of state-of-the-art scientific libraries such as the DUNE \cite{dune} library to perform the computations in an efficient way. Therefore, it leverages powerful Krylov subspace preconditioners and solvers that are generally used to solve large-scale and sparse linear systems because such systems often take a lot of time to solve; for example, they may even take 90\% or more of the simulation time for certain models. Efficient iterative solvers available in the DUNE library include GMRES and BiCGSTAB. To improve the convergence of an iterative solver, preconditioners are also used, with examples including AMG, domain decomposition, or multi-stage preconditioners \cite{saad1,saad2}, and the well-known ILU0, which is the default one used in OPM Flow. The choice for ILU as the default is due to the better performance, even against industry reference simulators, achieved by avoiding non-contributing computations \cite{ILU0-opt}. The Constrained Pressure Residual (CPR) \cite{CPR} preconditioner is a special multi-stage preconditioner, only used for reservoir simulation. It solves the mostly elliptic pressure system using an AMG preconditioner. To support large-scale processing, the DUNE library supports partitioning and execution on multiple cores and CPU nodes using OpenMP and MPI. However, it does not support GPU acceleration, and considering the adoption of GPUs in recent years to speed up many compute-intensive tasks, the question that we ask ourselves is if OPM would be able to efficiently leverage such technology as well and execute faster if the iterative solvers would be run on such an accelerator. 

GPUs have many small cores that can work together using the Same Instruction Multiple Data (SIMD) model. Each core performs the same instruction on different data. This allows the GPU to process massive amounts of data much faster than a CPU would, but only if the problem has enough parallelism. GPUs have shown large speedups in many different fields, like weather simulation\cite{gpu_weather}, FEM-based structural analysis\cite{gpu_struct}, numerically integrating ordinary differential equations (ODEs)\cite{gpu_ode}, chip design\cite{gpu_chip}, or machine learning with Tensorflow\cite{gpu_tensorflow}. However, although GPUs have proven very useful for several other application domains, it is not clear if they are useful for the reservoir simulation domain considering the specific sparsity. Specifically, the different sparsity patterns and the dependencies between reservoir model cells cause code divergence, which essentially means that cells placed into independent partitions, due to dependencies, are still executed sequentially. To analyze how big this problem is, we develop, integrate, and perform a thorough evaluation with both manual OpenCL kernels and various third-party libraries targeted at different GPU hardware to accelerate the linear solver of OPM Flow\cite{flow_paper} on a GPU. We target the default OPM ILU0 preconditioned BiCGStab linear solver used in OPM Flow. Concretely, the novelty of the paper is summarized as follows:

\begin{itemize}
    \item Develop open-source manual OpenCL and CUDA kernels for ILU0 preconditioned BICGStab solver.
    \item Develop a custom open-source bridge to integrate the custom-developed solvers as well as third-party libraries into the OPM Flow simulator. 
    \item Integrate both manual solvers and third-party libraries into OPM Flow, including amgcl, cuSparse, and rocSparse.
    \item Perform a thorough evaluation and comparison of all the developed and integrated libraries using three reservoir models with different sizes and running on GPU hardware from both Nvidia and AMD vendors. This is, to the best of our knowledge, the first work to present complete GPU results running on the open-source OPM Flow reservoir simulator using real-world use cases. 
\end{itemize}

The paper is organized as follows. First, in section \ref{sec:background}, we describe the background required to understand the OPM Flow simulator and the GPU preconditioners and iterative solver developed in this work. Then, in section \ref{sec:relres}, we highlight other interesting reservoir simulators and indicate what target processors they support. Section \ref{sec:implementation} describes the implementation details of how the different manual and third-party libraries were developed and integrated into OPM, including particularities specific to reservoir simulation, such as well contributions to the final solution. Section \ref{sec:experimentalresults} presents the experimental results and discusses the obtained performance numbers for the different libraries and GPU hardware tested. Finally, section \ref{sec:conclusion} summarizes the paper and lists future work. 
